This section briefly provides two additional results that the framework delivers, but defers their full derivations and proofs to the appendix.
Section~\ref{sec:extensions.aggregation} shows that, under additional assumptions, the firms' individual Leontief production functions can be aggregated into a single CES production function at the economy level.
This connects our firm-level, sequential model of production to models that study the implications of automation from a macro perspective.
Section~\ref{sec:extensions.computation} discusses how the cost minimization problems we have posed for the firm's short- and long-run optimization can actually be solved.


\subsection{CES Representation}
\label{sec:extensions.aggregation}

So far we have focused on a single firm's production: which steps it chains together, where it draws the boundaries of jobs, and what those jobs pay workers.
These adjustments do not stay inside the firm.
Repeated across firms, they are what determines how much labor the economy employs, how output responds to capital and to labor, and what becomes of both as AI quality improves.
These are the terms the automation literature works in, writing production functions directly at the level of the aggregate economy \citep{acemoglu2018, acemoglu2022artificial, acemoglu2024tasks}.

Our framework can be expressed in the same terms.
We distinguish two types of labor: \emph{manual labor}, which performs the manual portion of work, and \emph{AI management labor}, by which we mean the labor that verifies the AI's output at the end of a chain.\footnote{
The distinction between these two types of labor is deliberate rather than necessary.
One could carry the aggregation through with a single type of labor, pooling the two requirements into one input.
We keep them apart because the model already distinguishes the skill a step demands when a worker performs it, $\manualSkill{i}$, from the skill it demands when a worker verifies the output of an AI attempt at it, $\AIskill{i}$, so two types of labor are the natural counterpart of a distinction we have already drawn.
}
Appendix~\ref{sec:aggregation} shows that the firm-level technology developed earlier, in which the content of tasks and of jobs is chosen rather than given, can be aggregated to
\begin{align}
\label{eq:lr_macro_ces}
Y = \left(\theta_{\AIletter}\,\aggregateAIlabor^{\rho} + \theta_{\manualLetter}\,\aggregateManualLabor^{\rho} + (1-\theta_{\AIletter}-\theta_{\manualLetter})\,K^{\rho}\right)^{\frac{1}{\rho}},
\end{align}
over economy-wide AI management labor $\aggregateAIlabor$, manual labor $\aggregateManualLabor$, and capital $K$, with $Y$ denoting aggregate output.
This is the familiar constant elasticity of substitution (CES) form, with which the macroeconomic implications of automation are usually studied.

Arriving at Equation~\eqref{eq:lr_macro_ces} is not straightforward, and it is worth stating what makes it work.
No individual firm in our model substitutes AI management labor for manual labor, since producing a unit of output requires completing production steps in fixed proportions, so each firm's technology is Leontief.
Aggregate substitution comes instead from composition, in the manner of \citet{houthakker1955pareto} and \citet{levhari1968note}: firms have access to the same AI technology but differ in how effectively they are able to deploy it, and when AI management labor becomes relatively cheaper the economy's input mix moves because the set of firms for which using AI is worthwhile has moved, not because any firm has changed its own mix.
Appendix~\ref{sec:aggregation} obtains that dispersion from the order in which a firm has to commit, settling on its AI strategy and job design before it learns how effectively it can put AI to use, and derives in closed form the distribution of realized AI effectiveness under which Equation~\eqref{eq:lr_macro_ces} holds.\footnote{
The aggregate representation asks more of the economy than the firm-level argument does: capital productivity is common across firms, every firm organizes identically, realized effective AI quality is the single dimension along which firms differ, and the CES weights are tied to the skill-adjusted time requirements implied by the firm's chosen AI strategy and job design. See Appendix~\ref{sec:aggregation} for the details.
}

The micro and macro representations of production are therefore not in competition.
Our account of how chaining redraws tasks and jobs is not an alternative to the aggregate view of automation but a foundation for it, and the implications people draw from improving AI quality in a CES economy largely carry over.
The organization of work, however, does not disappear in aggregation; it is absorbed into the parameters of the CES function.
Two economies with identical production steps, arranged differently within workflows, aggregate to the same CES form in \eqref{eq:lr_macro_ces} but with different production parameters and input requirements.
Aggregate data therefore cannot reveal workflow structure, so objects such as the fragmentation index must be measured where the arrangement of work is still visible, which is what our empirical exercise in Section~\ref{sec:empirics} does.


\subsection{Computation of Firm's Short- and Long-run Optimums}
\label{sec:extensions.computation}

Throughout the paper we have assumed the firm can obtain its cost-minimizing arrangement without discussing how it is found.
This is not computationally trivial, as the firm cannot solve its optimization problem step by step due to chaining. 
Rather, it must consider all feasible arrangements, of which there are exponentially many.

That exponential count overstates the difficulty, however, because many of the configurations overlap heavily.
Any two arrangements that treat a fixed stretch of the workflow alike cost the same over that stretch, so the firm needs to calculate its cost just once.
Both the short- and long-run problems can therefore be formulated recursively and solved by dynamic programming, building the workflow up from the beginning one step at a time and reusing the work spent on an early stretch for every arrangement of what follows.
The two optimization problems differ only in how much a settled stretch has to remember about itself.
In the short run the configuration's time cost is all that matters, and the optimum can be calculated exactly.
In the long run the configuration's cost carries not just the time but also the skill of the worker whose job is still not terminated, and exactness gives way to an approximation the firm can make as tight as it likes.
We take the cost minimization of the two horizons in turn, collecting the proofs in Appendix~\ref{app:omitted_proofs}.

\paragraph{The Short Run Optimization.}
With the job and its wage fixed, the firm minimizes total expected time in Problem~\eqref{eq:totalcost_shortterm}.
For $k \leq m$, let $C[k]$ denote the minimum time needed to complete a hypothetical workflow containing only steps $1$ through $k$, so that $C[0] = 0$ and $C[m]$ is the optimum.
Step $k$ is either performed manually, or augmented as the endpoint of an AI chain reaching back to some earlier step $\ell$; in either case the steps before that point are optimized separately, which gives
\begin{align}
\label{eq:shortrun_recursion}
C[k] = \min\Biggl\{ \underbrace{C[k-1] + \manualTime{k}}_{\text{step $k$ manual}}, \ \underbrace{\min_{\ell < k}\ \Bigl( C[\ell] + \frac{\AItime{k}}{\prod_{i=\ell+1}^{k} q_i} \Bigr)}_{\text{step $k$ ends a chain begun at } \ell+1} \Biggr\}.
\end{align}
Because time costs enter Problem~\eqref{eq:totalcost_shortterm} additively, a stretch's time cost is all that the recursion in \eqref{eq:shortrun_recursion} needs from it.
Each of the $m$ values costs $O(m)$ to evaluate, since the chain ending at step $k$ may begin anywhere before it.
This makes the computational complexity of the algorithm $O(m^2)$.
Proposition~\ref{prop:shortrun_dp} formalizes this.
\begin{proposition}[Short-run AI Deployment]
\label{prop:shortrun_dp}
Given a workflow with $m$ steps, the AI strategy that solves Problem~\eqref{eq:totalcost_shortterm} can be calculated in time $O(m^2)$ via dynamic programming.
\end{proposition}

\paragraph{The Long Run Optimization.}
The firm now chooses job boundaries as well, and the objective in \eqref{eq:totalcost_with_handoff} multiplies each job's total skill by its total time.
A task's contribution to cost therefore depends on how much skill and how much time its job already carries, so the position alone no longer summarizes a stretch and we carry two further coordinates.
Let $R(i,\, \skillcost{},\, \timecost{})$ be the minimum cost of completing steps $1$ through $i$ when the worker currently accumulating tasks already holds total skill $\skillcost{}$ and total time $\timecost{}$, the latter including their hand-off, and let $V(i) = R(i,\, 0,\, \handofftime{i})$ be the corresponding cost when that worker's job closes at step $i$.
The base case is $R(0,\, \skillcost{},\, \timecost{}) = \skillcost{}\,\timecost{}$, and the firm faces three options at each step rather than two,
\begin{align}
\label{eq:longrun_recursion}
R(i,\, \skillcost{},\, \timecost{}) = \min\Biggl\{ \underbrace{\skillcost{}\,\timecost{} + V(i)}_{\text{close the job}}, \ \underbrace{R\bigl(i-1,\, \skillcost{} + \manualSkill{i},\, \timecost{} + \manualTime{i}\bigr)}_{\text{step $i$ manual}}, \ \underbrace{\min_{0 \le r < i}\ R\Bigl(r,\, \skillcost{} + \AIskill{i},\, \timecost{} + \tfrac{\AItime{i}}{\prod_{i'=r+1}^{i} q_{i'}}\Bigr)}_{\text{step $i$ ends a chain begun at } r+1} \Biggr\},
\end{align}
with the optimum at $V(m) = R(m,\, 0,\, 0)$.
Closing a job is where the accumulated skill and time are finally multiplied into a wage bill, which is why the recursion has to carry both alongside the position $i$.
That is also what makes the long run more complicated to solve.

Skill and time costs are continuous, so the recursion must distinguish between infinitely many situations rather than the $m$ of the short run, and it cannot be carried out to completion.
We therefore round each of the two costs to a power of $(1+\epsilon)$, which leaves finitely many of them to keep track of at the price of inflating the cost by at most a factor $(1+\epsilon)$.
Proposition~\ref{prop:totalcost_optimization_dp} formalizes the resulting approximation of the optimal solution.
\begin{proposition}[Long-run AI Deployment and Job Design]
\label{prop:totalcost_optimization_dp}
Fix a workflow with $m$ steps whose skill and time costs, and whose non-zero hand-off costs, lie in $[1/B, B]$ for some $B > 0$.
For any $\epsilon > 0$, a pair of AI strategy $\T$ and job design $\J$ whose long-run cost in \eqref{eq:totalcost_with_handoff} is within a factor $(1+\epsilon)$ of the minimum can be computed in time $O(m^2 \epsilon^{-2} \log^2(mB))$ via dynamic programming.
\end{proposition}

Together these results also suggest a use for the framework beyond what we do here.
Paired with detailed measurement of the time and skill content of individual tasks, the algorithms above could be run on real workflows to ask how a given advance in AI capability would redraw them.
